\documentclass[10pt,a4paper]{amsart}
\usepackage[margin=19mm]{geometry}
\usepackage{amsmath,amssymb,mathtools}
\usepackage[scr=rsfs,cal=euler]{mathalfa}
\usepackage{xfrac}
\usepackage[unicode,hidelinks,bookmarks=false]{hyperref}
\newcommand{\ie}{i.e.\ }
\newcommand{\cf}{cf.\ }
\newcommand{\resp}{respectively\ }
\newcommand{\Subsection}{\subsection}
\providecommand{\nobreakdash}{\nobreak\hbox{-}}
\setlength{\emergencystretch}{2em}
\multlinegap0pt
\usepackage{tikz}
\usepackage{amssymb}
\usepackage{enumerate}
\usepackage[shortlabels]{enumitem}
\usepackage[normalem]{ulem}
\newtheorem{proposition}{Proposition}
\newtheorem{lemma}{Lemma}
\newtheorem{theorem}{Theorem}
\newcommand{\einf}{\operatorname*{ess\,inf}}
\newcommand{\esup}{\operatorname*{ess\,sup}}
\title{\large{Energy inequalities for cutoff functions of $p$-energies on metric measure spaces}}
\author{Meng Yang}
\date{Source: arXiv:2507.08577v5 (CC BY 4.0)}
\begin{document}
\maketitle

\noindent\emph{Source and licence.} \large{Energy inequalities for cutoff functions of $p$-energies on metric measure spaces}. Version arXiv:2507.08577v5 (CC BY 4.0), distributed by arXiv under the Creative Commons Attribution 4.0 International licence, http://creativecommons.org/licenses/by/4.0/, as recorded in the arXiv licence metadata for this exact version and shown on the abstract page https://arxiv.org/abs/2507.08577v5. Full licence notice: \texttt{licenses/arxiv-2507.08577-CC-BY-4.0.txt}.\par\smallskip

\noindent\textit{Editorial scope.} Counted unit: the article's theorem thm\_Wolff (source0.tex lines 583--588). The article's complete original proof of it is delivered with it (source0.tex lines 650--709, one \texttt{proof} environment of 60 lines, which the article places away from the statement). No mathematics is written, repaired or re-derived: the statement and the proof are the article's own lines, in the article's own notation, and no retained line is substituted or re-flowed.  Retained uncounted as the source-local context the proof needs: lines 174--176; lines 284--286; lines 291--293; lines 296--298; lines 418--420; lines 438--446; lines 468--470; lines 565--567; lines 601--603; lines 607--609.  Nothing was omitted from any retained span, so the package carries no omission record.  No retained line contains a citation, so the wrapper carries no bibliography and no bibliography entry.  No retained line contains a figure environment, an image-inclusion command or drawing code, so no image is delivered and the package declares no asset dependency.  Editorial formatting only, in addition: an AMS \texttt{amsart} preamble that restates the article's own declarations and macros used by the retained lines, namely the environments \texttt{proposition}, \texttt{lemma}, \texttt{theorem} on separate counters, together with the standard packages those lines need.  The retained environments print in this document as Proposition 1, Lemma 1, Theorem 1 in order of appearance, whereas the article prints the same items with the numbering of its own full document.  The complete native source of the article as distributed in the arXiv e-print for this exact version is bundled unchanged as \texttt{sources/181422/source0.tex}.
\begin{equation*}\label{eq_VD}\tag*{VD}
V(x,2r)\le C_{VD}V(x,r)\text{ for any }x\in X,r>0.
\end{equation*}

\begin{equation*}\label{eq_BSP}\tag*{\text{BSP}}
\lambda_1(U)=\inf \left\{\frac{\mathcal{E}(u)}{\lVert u\rVert_{L^p(X;m)}^p}:u\in \mathcal{F}(U)\backslash \{0\}\right\}>0.
\end{equation*}

\begin{equation*}\label{eq_EHI}\tag*{\text{EHI}}
\esup_{B(x,r)}u\le C_H\einf_{B(x,r)}u.
\end{equation*}

\begin{equation*}\label{eq_EHIann}\tag*{$\text{EHI}_{\text{ann}}$}
\esup_{B(x,A_2r)\backslash B(x,A_1r)}u\le C_H\einf_{B(x,A_2r)\backslash B(x,A_1r)}u.
\end{equation*}

\begin{proposition}\label{prop_super_meas}
Let $U$ be an open set and $u\in\mathcal{F}$ superharmonic in $U$. Then there exists a positive Radon measure $\mu[u]$ on $U$ such that $\mathcal{E}(u;v)=\int_Uv\mathrm{d}\mu[u]$ for any $v\in\mathcal{F}\cap C_c(U)$, denoted by $-\Delta_pu=\mu[u]$ in $U$. Moreover, let $u_1$, $u_2\in\mathcal{F}$ be superharmonic in $U$, and $W$, $V$ two open sets satisfying that $W\subseteq \overline{W}\subseteq V\subseteq U$ and $u_1|_{V\backslash\overline{W}}=u_2|_{V\backslash\overline{W}}$, then $\mu[u_1](V)=\mu[u_2](V)$, that is, if $u_1=u_2$ near $V$, then $\mu[u_1](V)=\mu[u_2](V)$.
\end{proposition}

\begin{proposition}[Comparison principle]\label{prop_comparison}
Assume \ref{eq_VD}, \ref{eq_BSP}. Let $\lambda\ge0$. Let $U$ be a bounded open set and $u,v\in\mathcal{F}$ satisfying that
$$-\Delta_pu+\lambda|u|^{p-2}u\ge-\Delta_pv+\lambda|v|^{p-2}v\text{ in }U,$$
that is,
$$\mathcal{E}(u;\varphi)+\lambda\int_X|u|^{p-2}u\varphi\mathrm{d} m\ge\mathcal{E}(v;\varphi)+\lambda\int_X|v|^{p-2}v\varphi\mathrm{d} m$$
for any non-negative $\varphi\in\mathcal{F}(U)$, and $(v-u)_+\in \mathcal{F}(U)$, or
$$\widetilde{u}\ge \widetilde{v}\text{ q.e. on }W\backslash U\text{ for some open set }W\supseteq \overline{U}\supseteq U.$$
Then $u\ge v$ in $U$.
\end{proposition}

\begin{proposition}\label{prop_Pois_modi}
Assume \ref{eq_VD}, \ref{eq_BSP}. Let $U$ be a bounded open set, $u\in\mathcal{F}$ superharmonic in $U$, and $V$ an open subset of $U$. Then $H^{V}u\in\mathcal{F}$ is superharmonic in $U$, harmonic in $V$, and $H^{V}u\le u$ in $U$.
\end{proposition}

\begin{lemma}\label{lem_FU_modi}
Assume \ref{eq_VD}, \ref{eq_BSP}. Let $U$ be a bounded open set and $u\in\mathcal{F}$ superharmonic in $U$. Then there exists $v\in\mathcal{F}(U)$ which is non-negative superharmonic in $U$ such that $\mu[v]=\mu[u]$ in $U$. Moreover, if $(-u)_+\in \mathcal{F}(U)$ or $\widetilde{u}\ge0$ q.e. on $W\backslash U$ for some open set $W\supseteq\overline{U}\supseteq U$, then $v\le u$ in $U$.
\end{lemma}

\begin{equation}\label{eq_Wolff1}
\einf_{\frac{1}{2}B}u\ge \frac{1}{C}\left(\frac{\mu[u]\left(\frac{9}{10}B\right)}{\mathrm{cap}\left(\frac{1}{2}B,X\backslash B\right)}\right)^{1/(p-1)}.
\end{equation}

\begin{equation}\label{eq_Wolff2}
\esup_{\frac{11}{20}B\backslash {\frac{9}{20}B}}u-\esup_{\frac{11}{10}B\backslash {\frac{9}{10}B}}u\le C\left(\frac{\mu[u](B)}{\mathrm{cap}\left(\frac{1}{2}B,X\backslash B\right)}\right)^{1/(p-1)}.
\end{equation}

\begin{theorem}\label{thm_Wolff}
Assume \ref{eq_VD}, \hyperlink{eq_LSC}{\text{LSC}}, \ref{eq_BSP}, \ref{eq_EHI}, \ref{eq_EHIann}, \hyperlink{eq_cap}{$\text{cap}(\Psi)$}. Then there exists $C>0$ such that for any bounded open set $U$, for any $u\in\mathcal{F}$ which is non-negative superharmonic in $U$, for any Lebesgue point $x_0\in U$ of $u$, for any $R>0$ with $B(x_0,4R)\subseteq U$, we have
\begin{equation*}\label{eq_Wolff}\tag*{$\text{Wol}$}
\frac{1}{C}\mathcal{W}^{\mu[u]}(x_0,R)\le u(x_0)\le C \left(\einf_{B(x_0,R)}u+\mathcal{W}^{\mu[u]}(x_0,2R)\right).
\end{equation*}
\end{theorem}

\begin{proof}[Proof of Theorem \ref{thm_Wolff}]
By \hyperlink{eq_LSC}{\text{LSC}}, we may assume that $u$ is lower semi-continuous in $U$. Firstly, we prove the lower bound. For any $r\in(0,2R]$, let $B=B(x_0,r)$, then $u$ is non-negative superharmonic in $2B$. By Lemma \ref{lem_FU_modi}, there exists $u_1\in \mathcal{F}(B)$ which is non-negative superharmonic in $B$ such that $\mu[u_1]=\mu[u]$ in $B$, moreover, $u_1\le u$ in $B$. Since $\widetilde{u}_1=0\le \widetilde{u}-\einf_{\frac{3}{2}B}u$ q.e. on $\frac{3}{2}B\backslash B$, by Proposition \ref{prop_comparison}, we have
$$\einf_{\frac{3}{4}B}u_1\le \einf_{\frac{3}{4}B}u-\einf_{\frac{3}{2}B}u.$$
Let $u_2=H^{\frac{7}{8}B}u_1$, then by Proposition \ref{prop_Pois_modi}, we have $u_2$ is non-negative superharmonic in $B$, harmonic in $\frac{7}{8}B$, $\widetilde{u}_2=\widetilde{u}_1$ q.e. on $B\backslash \frac{7}{8}B$, and $u_2\le u_1$ in $B$. By (\ref{eq_Wolff1}), we have
$$\einf_{\frac{1}{2}B}u_2\ge \frac{1}{C_1}\left(\frac{\mu[u_2]\left(\frac{9}{10}B\right)}{\mathrm{cap}\left(\frac{1}{2}B,X\backslash B\right)}\right)^{1/(p-1)},$$
where $C_1$ is the positive constant appearing therein. Since $u_2=u_1$ in $\frac{9}{10}B\backslash \overline{\frac{7}{8}B}$, by Proposition \ref{prop_super_meas}, we have
$$\mu[u_2]\left(\frac{9}{10}B\right)=\mu[u_1]\left(\frac{9}{10}B\right)=\mu[u]\left(\frac{9}{10}B\right)\ge \mu[u]\left(\frac{1}{2}B\right).$$
Since $u_2$ is non-negative harmonic in $\frac{7}{8}B$, by \ref{eq_VD}, \ref{eq_EHI}, there exists $C_2>0$ depending only on $C_{VD}, C_H, A_H$ such that
$$\einf_{\frac{1}{2}B}u_2\le \esup_{\frac{3}{4}B}u_2\le C_2\einf_{\frac{3}{4}B}u_2\le C_2\einf_{\frac{3}{4}B}u_1.$$
Hence
$$\einf_{\frac{3}{4}B}u-\einf_{\frac{3}{2}B}u\ge \frac{1}{C_1C_2}\left(\frac{\mu[u]\left(\frac{1}{2}B\right)}{\mathrm{cap}\left(\frac{1}{2}B,X\backslash B\right)}\right)^{1/(p-1)}.$$
By \ref{eq_VD}, \hyperlink{eq_cap}{$\text{cap}(\Psi)$}, there exists $C_3>0$ depending only on $C_{\Psi}, C_{VD}, A_{cap}, C_{cap}$ such that
$${\mathrm{cap}\left(\frac{1}{2}B,X\backslash B\right)}\le C_3{\mathrm{cap}\left(\frac{1}{4}B,X\backslash \frac{1}{2}B\right)},$$
hence
$$\einf_{\frac{3}{4}B}u-\einf_{\frac{3}{2}B}u\ge\frac{1}{C_1C_2C_3^{1/(p-1)}}\left(\frac{\mu[u]\left(\frac{1}{2}B\right)}{\mathrm{cap}\left(\frac{1}{4}B,X\backslash \frac{1}{2}B\right)}\right)^{1/(p-1)}.$$
Then for any $k=0,1,\ldots$, we have
\begin{align*}
&\einf_{B(x_0,\frac{3}{2^{k+1}}R)}u-\einf_{B(x_0,\frac{3}{2^{k}}R)}u\\
&\ge\frac{1}{C_1C_2C_3^{1/(p-1)}}\left(\frac{\mu[u]\left(B(x_0,\frac{1}{2^{k}}R)\right)}{\mathrm{cap}\left(B(x_0,\frac{1}{2^{k+1}}R),X\backslash B(x_0,\frac{1}{2^{k}}R)\right)}\right)^{1/(p-1)}.
\end{align*}
Summing over $k=0,1,\ldots,l$, we have
\begin{align*}
&\einf_{B(x_0,\frac{3}{2^{l+1}}R)}u\ge \einf_{B(x_0,\frac{3}{2^{l+1}}R)}u-\einf_{B(x_0,{3}R)}u\\
&\ge\frac{1}{C_1C_2C_3^{1/(p-1)}}\sum_{k=0}^l\left(\frac{\mu[u]\left(B(x_0,\frac{1}{2^{k}}R)\right)}{\mathrm{cap}\left(B(x_0,\frac{1}{2^{k+1}}R),X\backslash B(x_0,\frac{1}{2^{k}}R)\right)}\right)^{1/(p-1)}.
\end{align*}
Since $x_0$ is a Lebesgue point of $u$, we have
\begin{align*}
&u(x_0)=\lim_{l\to+\infty}\frac{1}{m({B(x_0,\frac{3}{2^{l+1}}R)})}\int_{{B(x_0,\frac{3}{2^{l+1}}R)}}u\mathrm{d} m\ge\varlimsup_{l\to+\infty}\einf_{B(x_0,\frac{3}{2^{l+1}}R)}u\\
&\ge\frac{1}{C_1C_2C_3^{1/(p-1)}}\sum_{k=0}^{+\infty}\left(\frac{\mu[u]\left(B(x_0,\frac{1}{2^{k}}R)\right)}{\mathrm{cap}\left(B(x_0,\frac{1}{2^{k+1}}R),X\backslash B(x_0,\frac{1}{2^{k}}R)\right)}\right)^{1/(p-1)}.
\end{align*}

Secondly, we prove the upper bound. Let $\omega=\cup_{k=0}^{+\infty}(B(x_0,\frac{5}{4}\frac{1}{2^k}R)\backslash\overline{B(x_0,\frac{3}{4}\frac{1}{2^k}R)})$ and $v=H^\omega u$, then by Proposition \ref{prop_Pois_modi}, we have $v$ is non-negative superharmonic in $B(x_0,4R)$, harmonic in $\omega$, $\widetilde{v}=\widetilde{u}$ q.e. on $B(x_0,4R)\backslash\omega$, and $v\le u$ in $B(x_0,4R)$. By Proposition \ref{prop_super_meas}, for any $r\in(0,{R}]$, we have $\mu[v](B(x_0,r))\le\mu[u](B(x_0,2r))$. By (\ref{eq_Wolff2}), for any $k=0,1,\ldots$,
\begin{align*}
&\esup_{B(x_0,\frac{11}{10}\frac{1}{2^{k+1}}R)\backslash {B(x_0,\frac{9}{10}\frac{1}{2^{k+1}}R)}}v\\
&\le\esup_{B(x_0,\frac{11}{10}\frac{1}{2^{k}}R)\backslash {B(x_0,\frac{9}{10}\frac{1}{2^{k}}R)}}v+C_1 \left(\frac{\mu[v](B(x_0,\frac{1}{2^k}R))}{\mathrm{cap}\left(B(x_0,\frac{1}{2^{k+1}}R),X\backslash B(x_0,\frac{1}{2^k}R)\right)}\right)^{1/(p-1)}\\
&\le\esup_{B(x_0,\frac{11}{10}\frac{1}{2^{k}}R)\backslash {B(x_0,\frac{9}{10}\frac{1}{2^{k}}R)}}v+C_1 \left(\frac{\mu[u](B(x_0,\frac{1}{2^{k-1}}R))}{\mathrm{cap}\left(B(x_0,\frac{1}{2^{k+1}}R),X\backslash B(x_0,\frac{1}{2^k}R)\right)}\right)^{1/(p-1)}.
\end{align*}
Summing over $k=0,1,\ldots,l$, we have
\begin{align*}
&\esup_{B(x_0,\frac{11}{10}\frac{1}{2^{l+1}}R)\backslash {B(x_0,\frac{9}{10}\frac{1}{2^{l+1}}R)}}v\\
&\le\esup_{B(x_0,\frac{11}{10}R)\backslash {B(x_0,\frac{9}{10}R)}}v+C_1\sum_{k=0}^l\left(\frac{\mu[u](B(x_0,\frac{1}{2^{k-1}}R))}{\mathrm{cap}\left(B(x_0,\frac{1}{2^{k+1}}R),X\backslash B(x_0,\frac{1}{2^k}R)\right)}\right)^{1/(p-1)}.
\end{align*}
Since $v$ is non-negative harmonic in $B(x_0,\frac{5}{4}R)\backslash \overline{B(x_0,\frac{3}{4}R)}$, by \ref{eq_EHIann}, there exists $C_4>0$ such that
\begin{equation*}
\esup_{B(x_0,\frac{11}{10}R)\backslash {B(x_0,\frac{9}{10}R)}}v\le C_4\einf_{B(x_0,\frac{11}{10}R)\backslash {B(x_0,\frac{9}{10}R)}}v\le C_4\einf_{B(x_0,\frac{11}{10}R)\backslash {B(x_0,R)}}v,
\end{equation*}
where
$$\einf_{B(x_0,\frac{11}{10}R)\backslash {B(x_0,R)}}v\le\einf_{B(x_0,R)}v\le\einf_{B(x_0,R)}u.$$
By the lower semi-continuity, we have $u(x_0)\le\lim_{l\to+\infty}\einf_{B(x_0,\frac{9}{10}\frac{1}{2^{l}}R)}u$, where
\begin{align*}
&\einf_{B(x_0,\frac{9}{10}\frac{1}{2^{l}}R)}u= \einf_{B(x_0,\frac{9}{10}\frac{1}{2^{l}}R)\backslash B(x_0,\frac{11}{10}\frac{1}{2^{l+1}}R)}u=\einf_{B(x_0,\frac{9}{10}\frac{1}{2^{l}}R)\backslash B(x_0,\frac{11}{10}\frac{1}{2^{l+1}}R)}v\\
&\le \einf_{B(x_0,\frac{11}{10}\frac{1}{2^{l+1}}R)\backslash B(x_0,\frac{9}{10}\frac{1}{2^{l+1}}R)}v\le \esup_{B(x_0,\frac{11}{10}\frac{1}{2^{l+1}}R)\backslash B(x_0,\frac{9}{10}\frac{1}{2^{l+1}}R)}v,
\end{align*}
hence
$$u(x_0)\le C_4\einf_{B(x_0,R)}u+C_1\sum_{k=0}^{+\infty}\left(\frac{\mu[u](B(x_0,\frac{1}{2^{k-1}}R))}{\mathrm{cap}\left(B(x_0,\frac{1}{2^{k+1}}R),X\backslash B(x_0,\frac{1}{2^k}R)\right)}\right)^{1/(p-1)}.$$
By \ref{eq_VD}, \hyperlink{eq_cap}{$\text{cap}(\Psi)$}, there exists $C_5>0$ depending only on $C_\Psi,C_{VD},C_{cap},A_{cap}$ such that for any $k=0,1,\ldots$, we have
$$\mathrm{cap}\left(B(x_0,\frac{1}{2^{k+1}}R),X\backslash B(x_0,\frac{1}{2^k}R)\right)\ge \frac{1}{C_5}\mathrm{cap}\left(B(x_0,\frac{1}{2^{k}}R),X\backslash B(x_0,\frac{1}{2^{k-1}}R)\right).$$
Hence
$$u(x_0)\le C_4\einf_{B(x_0,R)}u+C_1C_5^{1/(p-1)}\sum_{k=0}^{+\infty}\left(\frac{\mu[u](B(x_0,\frac{1}{2^{k}}2R))}{\mathrm{cap}\left(B(x_0,\frac{1}{2^{k+1}}2R),X\backslash B(x_0,\frac{1}{2^k}2R)\right)}\right)^{1/(p-1)}.$$
\end{proof}

\end{document}
